% byu-coi-briefing.tex
%
% THE BRIEFING. Background for a department chair, a dean, and the AAVP for
% Research. This document is NOT submitted and is NOT signed. It carries the
% approval sequence, the signing tiers, the reasoning behind the plan's
% conditions, and the questions still open.
%
% The instrument itself is byu-coi-plan.tex. Anything that reads as an
% obligation the Faculty Member is taking on belongs there, not here.
%
% No em or en dashes (see the house rule in vcl-doc.sty).
% Compile with pdflatex, THREE times, then copy the PDF to the repo root.
%
\documentclass[11pt]{article}
\usepackage{vcl-doc}

\hypersetup{
  pdftitle={Getting Approval at BYU: background for the management plan},
  pdfsubject={Approval sequence and background for a faculty co-founder conflict-of-interest plan},
  pdfkeywords={BYU, conflict of interest, faculty founder, technology transfer, approvals}
}
\vclrunninghead{Getting Approval at BYU, Vertical Cloud Lab}

\begin{document}

% ================================================================== title ==
\thispagestyle{plain}

{\color{byunavy}\hrule height 2.4pt}
\vspace{12pt}
{\LARGE\bfseries\color{byunavy} Getting Approval at BYU}\\[6pt]
{\normalsize Background for the conflict-of-interest management plan}
\vspace{8pt}

{\color{byunavy}\hrule height 0.8pt}
\vspace{9pt}

A faculty member who co-founds a venture-funded company while remaining a
full-time professor needs a defined set of university approvals. This
document lays out that path: what happens in what order, who has authority to
sign, what the university might gain, and what is still unsettled.

\noticebox{\textbf{This is background, not the plan.} The conflict-of-interest
management plan is a separate six-page document, approved and signed on its
own. BYU's \href{https://spo.byu.edu/policies}{Financial Disclosure Form} asks
whether ``a proposed conflict of interest management plan is attached,'' and
that plan is the attachment. Nothing in this briefing is submitted with it.
The
\href{https://policy.byu.edu/view/financial-conflict-of-interest-in-research-policy}{Financial
Conflict of Interest in Research Policy} has the faculty member develop that
plan in consultation with the Research Administration Office for approval by
the Associate Academic Vice President for Research (``the AAVP''). BYU
publishes no plan template, so the plan is organized on the headings policy
requires the university to be able to report, with the reasoning for one of
them (how the conditions safeguard objectivity) kept here instead, checked
against the formats peer universities publish
(\href{https://provost.uiowa.edu/management-plan-conflict-commitment-andor-conflict-interest-workplace}{Iowa},
\href{https://www.research.colostate.edu/integrity-security-compliance/wp-content/uploads/sites/25/2026/03/COI-COC-management-plan_2025.v2.docx}{Colorado
State},
\href{https://research.unl.edu/researchcompliance/wp-content/uploads/sites/22/2024/06/FOR-WEB-Mgmt-Plan-Template_FINAL.pdf}{Nebraska}).}

\subhead{Offices named here}

{\small
\setlength{\extrarowheight}{2pt}
\begin{tabular}{@{} >{\RaggedRight\arraybackslash}p{2.1in} >{\RaggedRight\arraybackslash}p{4.35in} @{}}
\textbf{Research Administration Office} &
negotiates and signs research agreements on the university's behalf; first
reviewer of financial disclosures. \\
\textbf{Technology Transfer Office} &
licenses BYU-owned inventions; reviews sponsored research and consulting
agreements that touch intellectual property. \\
\textbf{AAVP for Research} &
decides whether a disclosed interest is a conflict, approves every management
plan, and makes the final decision on any transfer of BYU technology. \\
\textbf{Regulatory Accounting} &
approves cost-recovery accounts and rates; one of three offices that must
approve an employee acting as a vendor to BYU. \\
\end{tabular}}

% ================================================== what policy already fixes
\parthead{What BYU policy already settles}
\label{sec:numbers}

Two numbers decide most of this before anyone exercises judgment: the faculty
member holds \textbf{more than ten percent} of the company, and the company
has raised \textbf{more than \$3 million}. Four consequences follow, and the plan
accepts all of them rather than seeking exceptions.

\begin{enumerate}[leftmargin=1.9em]
\item \textbf{The faculty member stays out of every deal between BYU and the
Company.} Above ten percent the company is a ``Related Party'' under the
\href{https://policy.byu.edu/view/conflict-of-time-or-interest-policy}{Conflict
of Time or Interest Policy}, and ``negotiating, influencing, or attempting to
influence the negotiations of contracts or agreements between the university
and a third party'' for its benefit is on that policy's \emph{prohibited}
list, the conflicts BYU says cannot be managed at all. Someone else negotiates
and signs on each side.
\item \textbf{Advisory roles, not management roles.} A faculty member not on
leave is ``not to be line officers in outside companies''
(\href{https://finserve.byu.edu/00000193-6566-d995-a9db-eff6a47a0001/sponsored-programs-handbook}{Sponsored
Projects Handbook}), the same line
\href{https://doresearch.stanford.edu/policies/research-policy-handbook/conflicts-commitment-and-interest/faculty-policy-conflict-commitment-and-interest}{Stanford
draws}: advisory roles yes, management roles and titles no.
\item \textbf{No voting board seat.} The
\href{https://policy.byu.edu/view/intellectual-property-policy}{Intellectual
Property Policy} allows a voting seat on the board of a company that licenses
BYU technology only if the company ``has not raised more than \$3,000,000''
and has no non-accredited investors, and even then only ``on rare occasions''
with prior written approval from the dean, the AAVP, \emph{and} the Technology
Transfer Office. This company is past that ceiling. The policy says nothing
about non-voting observer seats, which means such a seat is not authorized by
default either; if one is wanted it has to be granted in writing.
\item \textbf{The inventor's royalty share is given up.} Under the same
policy, ``a principal equity owner is one who controls ten percent or more
equity in the company,'' and where a developer becomes one, ``the developer
will forfeit his or her distribution.'' That is the personal 45~percent share
of BYU's licensing income, and with it the option to steer that share, with
university matching, into a research account. The policy sets out no way to
apply for an exception. The faculty member's financial interest is the company
stock alone.
\end{enumerate}

\noindent \textbf{One thing policy does not settle by prohibition: travel.} An
advisory role means occasional trips to the company, its partners, and its
customers, and a rule requiring the faculty member to absorb the airfare
personally would be either broken or quietly ignored. BYU treats sponsor-paid
travel as something to report. The
\href{https://spo.byu.edu/policies}{Financial Disclosure Form} gives it an item
of its own, asking whether ``travel expenditures \ldots{} were paid on your
behalf or reimbursed to you by a sponsor, vendor, or subcontractor of the
sponsored program activity,'' and then collecting the sponsor, purpose, dates,
and destination of each trip. The accompanying instructions class such travel
as a significant financial interest, which is the category BYU manages through
disclosure and a plan rather than by banning it, and they exempt ``travel paid
for by BYU or sponsored research funds managed by BYU'' from disclosure
altogether. So the plan bounds travel rather than forbidding it: documented
actual cost only, the days counted against the four-day allowance, no trip paid
twice, and university business kept on university funds (clause 1.4). Anything
beyond the trip itself stays within the Business Gifts and Entertainment
Policy's ``nominal value,'' which admits no cash or cash equivalents. The
distinction that does the work is between reimbursement, which leaves the
faculty member no better off than before the trip, and compensation, which is
what clause 1.3 rules out.

% ================================================================== steps ==
\parthead{The approval sequence}

``Department approval'' comes in layers rather than as a single event. The
department controls the disclosure, the written authorization of the outside
role, and the management plan; each later transaction adds university-level
approvers. Two rules run through all of it: \textbf{disclose before acting},
and \textbf{every arrangement between BYU and the Company lives in a written
agreement signed by the university}, never by the faculty member.

{\footnotesize
\setlength{\extrarowheight}{2pt}
\begin{longtable}{@{} >{\bfseries}p{0.22in} >{\RaggedRight\arraybackslash}p{2.75in} >{\RaggedRight\arraybackslash}p{1.95in} >{\RaggedRight\arraybackslash}p{1.15in} @{}}
\toprule
{\color{byunavy}\scshape\#} &
{\color{byunavy}\scshape What happens} &
{\color{byunavy}\scshape Who approves} &
{\color{byunavy}\scshape When} \\
\midrule
\endfirsthead
\toprule
{\color{byunavy}\scshape\#} &
{\color{byunavy}\scshape What happens} &
{\color{byunavy}\scshape Who approves} &
{\color{byunavy}\scshape When} \\
\midrule
\endhead
\midrule
\multicolumn{4}{r@{}}{\itshape continued on next page}\\
\endfoot
\bottomrule
\endlastfoot
1 & Informal briefing of the department chair, with this document as the
agenda & none & first \\
\addlinespace
2 & Outside-activity disclosure filed in
\href{https://facultyprofile.byu.edu/}{Faculty Profile} (the ``Conflict of
Interest and Conflict of Time'' screen, Annual Review section) & filed, not
approved & \textbf{before} the role begins; renewed annually \\
\addlinespace
3 & Written authorization of the outside role (operating a business; advisory
service) & department chair (policy allows chair, dean, or supervisor) &
before taking the role \\
\addlinespace
4 & The plan as the \textbf{outside-activity} conflict plan, written with the
chair, who submits it to the dean & dean & together with step~3 \\
\addlinespace
5 & \href{https://spo.byu.edu/policies}{Financial Disclosure Form}, with the
plan attached as the \textbf{research} conflict plan & chair, then the dean's
office, then the Research Administration Office; \textbf{the AAVP approves} &
before any research proposal involving the company; updates within 30~days \\
\addlinespace
6 & Conflict-of-interest training: BYU's online
\href{https://spo.byu.edu/responsible-conduct-of-research}{research-conduct
course}, required of all faculty involved in sponsored research & none &
refreshed every four years, and immediately on joining research, on a policy
change, or after a compliance finding \\
\addlinespace
7 & License or assignment of BYU technology & chair, dean, Technology Transfer
Office, and General Counsel review; \textbf{the AAVP decides} & at deal time \\
\addlinespace
8 & Company-sponsored research agreement & Research Administration Office
negotiates and signs; chair, dean, and the AAVP approve the proposal & when
company money funds lab work \\
\addlinespace
9 & Company equipment placed in the lab & Research Administration Office signs
a loan agreement or research-agreement term & before delivery \\
\addlinespace
10 & Lab sells instrument time at cost-recovery rates
(\href{https://spo.byu.edu/00000190-3687-d777-a591-37af392d0001/rechargecenterpolicyprocedures-pdf}{Recharge
Center Policy}) & chair, then dean, then Regulatory Accounting (rates
re-approved annually) & before billing outside users \\
\addlinespace
11 & The company sells anything to BYU & employee-vendor approval by
Purchasing \& Travel, Compensation, and Regulatory Accounting & before any BYU
purchase \\
\end{longtable}}

\begin{itemize}
\item \textbf{The standard the chair or dean applies is written down.} They
``do not approve activities they believe would unduly limit a faculty member's
availability to students or compromise\ldots\ teaching, scholarship, or
citizenship performance''
(\href{https://policy.byu.edu/view/conflict-of-time-or-interest-policy}{Conflict
of Time or Interest Policy}). The step-1 briefing should therefore lead with
the time budget and the teaching plan, not the company.
\item \textbf{Disclosure comes before the activity.} A plan approved early,
while the stakes are small, is the dated record that answers scrutiny later if
the company succeeds.
\item \textbf{Steps 7 through 11 recur with each transaction}, and the faculty
member is out of the negotiation every time.
\end{itemize}

% ================================================================ signing ==
\parthead{Who is allowed to sign}

Signature authority at BYU is set by what an agreement costs and how long it
binds the university, not by what kind of agreement it is. A faculty member
sits in the bottom row, which is why the plan's equipment and research
agreements are signed well above the lab.

{\footnotesize
\setlength{\extrarowheight}{2pt}
\begin{tabularx}{\textwidth}{@{} >{\RaggedRight\arraybackslash}X >{\RaggedRight\arraybackslash}p{1.1in} >{\RaggedRight\arraybackslash}p{1.45in} @{}}
\toprule
{\color{byunavy}\scshape Who signs} &
{\color{byunavy}\scshape Up to} &
{\color{byunavy}\scshape Can bind BYU for} \\
\midrule
President (or delegate) & no limit & no limit \\
President's Council member (or delegate) & \$500,000 & 5 years \\
Associate or Assistant Vice President, or a Dean or Director reporting to a
President's Council member & \$250,000 & 3 years \\
Any other employee, acting within the scope of employment and with the
supervisor's approval (\textbf{this is the faculty member}) & \$100,000 & 1
year \\
\bottomrule
\end{tabularx}}

\vspace{2pt}
\begin{itemize}
\item \textbf{A multi-year equipment placement is out of the lab's hands.}
Anything past one year needs a dean or director at minimum, and past three
years or \$250,000 it goes to a President's Council member.
\item \textbf{The agreement cannot be sliced to fit.} ``A Transaction may not
be divided between two or more Legal Documents to avoid authorization
limits,'' and ``a series of related Transactions with the same party for the
same subject matter is considered a single Transaction.'' Five one-year
renewals of the same placement is one five-year transaction.
\item \textbf{Certain provisions pull in other offices regardless of size.}
The Office of General Counsel must review anything that indemnifies another
party, sets non-Utah governing law, limits BYU's liability, gives a third
party access to non-public university data, or transfers a university
ownership interest in a company. That last one applies the moment BYU takes
Company stock. Risk Management must review any commitment to carry or show
insurance, any indemnity, and any activity carrying a real risk of injury or
loss. A department cannot buy its own coverage.
\end{itemize}

% =============================================================== benefits ==
\parthead{What the university might gain, and what each depends on}

None of these is a commitment. Each depends on a separate agreement or
approval the plan does not grant, and a benefit is not a reason to approve a
conflict.

{\footnotesize
\setlength{\extrarowheight}{2pt}
\begin{tabularx}{\textwidth}{@{} >{\RaggedRight\arraybackslash}p{2.2in} >{\RaggedRight\arraybackslash}X @{}}
\toprule
{\color{byunavy}\scshape Possible benefit} &
{\color{byunavy}\scshape What it depends on} \\
\midrule
Research students get hands-on time on instruments the university would find
hard to buy &
A university-signed equipment agreement fixing ownership, insurance,
maintenance, access, and removal. The Company can withdraw its equipment;
nothing here is permanent. \\
\addlinespace
Instrument time sold to outside users, and Company sponsorship, fund student
wages, equipment, and student travel &
A cost-recovery account approved by the chair, dean, and Regulatory
Accounting, rates re-approved annually, and the tax treatment of outside sales
settled before billing. Spending follows sponsor terms and university budget
rules, not the lab's preferences. \\
\addlinespace
New research funding through independently negotiated sponsored agreements &
Each agreement negotiated and signed by the Research Administration Office
with the faculty member recused. Any federal small-business subaward that
comes later is federal money and carries federal rules. \\
\addlinespace
Industry experience for students: internships, applied problems, jobs &
Participation stays voluntary, is reviewed by an independent advocate, and is
protected from retaliation. \\
\addlinespace
Student work stays publishable, because the Company's model is open source &
The lab's and the Company's published open-source policies, both adopted in
advance. Anything outside them still goes to the Technology Transfer Office
first, and export controls, third-party rights, and confidential information a
sponsor supplies apply regardless of the open-source default. \\
\bottomrule
\end{tabularx}}

% ============================================================== safeguards ==
\parthead{Why the plan's conditions are enough}

BYU has to be able to say how an approved plan protects the objectivity of the
research, and five answers carry that weight. \textbf{Influence over a deal
between BYU and the Company} is removed by total recusal on both sides (clauses
3.1, 3.2), which above ten percent is not a judgment call: policy lists
influencing such negotiations among the conflicts it says cannot be managed at
all. \textbf{Payment tied to Company outcomes} does not exist: no salary, fee,
bonus, or honorarium (1.3), a forfeited royalty share (1.5), and travel
reimbursed at documented cost, counted against the four-day allowance and
reported trip by trip (1.4). What is left is unsellable private stock, which no
single contract term turns into money. \textbf{Research shaped, suppressed, or
delayed} is checked by independent review before work begins and annually
after, with the chair able to pause it (4.2), by the bar on embargoes and on
confidentiality at a defense (5.2), and by disclosure on every publication
(4.3). \textbf{Pressure on students} is answered by real alternatives and a bar
on retaliation (5.2), an advocate outside the lab who meets them alone (5.1),
review of any Company offer (5.4), and one employer at a time (5.3).
\textbf{The two roles merging} is prevented by Company work on Company systems
(6.1), personnel and equipment admitted only under signed agreements (6.3,
6.4), and a Company officer signing the promises made about Company conduct
(8.3).

% =============================================================== formation ==
\parthead{Company formation documents, and the line they must respect}

Standard startup formation packages have founders buy their shares partly with
intellectual property, and have everyone who works for the company assign
inventions made \emph{before} as well as after incorporation. Stripe Atlas,
the most common of them, says both in its own
\href{https://docs.stripe.com/atlas/incorporation-documents}{document list}:
``Atlas founders pay for their shares with the intellectual property (IP) they
have developed on behalf of the company,'' and its invention-assignment form
reaches ``relevant IP created both before and after incorporation.''

A contribution like that is only as good as the founder's title to it, and
BYU's claim attaches automatically: employees ``agree to assign and do hereby
assign all right, title, and interest in any Intellectual Property, including
future Intellectual Property,'' reaching work ``related to the current or
demonstrably anticipated business, research, or development of the
university''
(\href{https://policy.byu.edu/view/intellectual-property-policy}{Intellectual
Property Policy}). So:

\begin{itemize}
\item Anything within that description was the university's before the Company
existed. It reaches the Company only through a license negotiated without the
faculty member, never by founder assignment.
\item \textbf{Founder shares are paid for in cash, and nothing is listed as
contributed.} The cash is already in the standard package. Atlas founders
``purchase their shares with a combination of money and IP''
(\href{https://stripe.com/guides/atlas/equity}{Stripe's own guide to founder
equity}), and the money side is ``cash equal to the purchase price of the
shares,'' which is what ``ensures that the shares are fully paid''
(\href{https://stripe.com/files/atlas/orrick-legal-guide.pdf}{the annotated
legal guide Stripe publishes with it}). At \$0.00001 a share that is a trivial
sum, and it means the assignment bundled into the purchase is not what the
shares rest on: contributing nothing still leaves them fully paid.
Contributing intellectual property instead would require valuing the very
thing whose ownership is in question.
\item What the faculty member may contribute is limited to work created
outside university employment: no BYU funding, facilities, equipment, systems,
personnel, or work time, and never disclosed to or licensed from the
university.
\item That boundary is settled through BYU's own release procedure before any
formation or invention-assignment document is signed, not asserted afterward.
The \href{https://techtransfer.byu.edu/resources}{invention disclosure form}
sets it out: an invention ``not within the field of employment'' that ``did
not employ BYU resources'' is not disclosed on the form; instead ``a brief,
non-confidential description of the Invention should be submitted to the
Technology Transfer Office with a copy to the employee's chairman or director
to request release of the Invention to the employee.'' ``Demonstrably
anticipated'' can otherwise reach personal-time work in the lab's own field,
and the university's view of that line is the one that governs.
\end{itemize}

Stripe flags the mismatch itself (``if you have any prior inventions you'd
like to exclude from being assigned to the company, Atlas might not be right
for you''), and supplies the fix without leaving the service: ``you can work
with them to edit the documents you receive from Stripe Atlas.'' The automated
stock tool is one of two routes; the other is the annotated templates issued on
the same dashboard for counsel to edit.

Note which document does the work. Paying cash and contributing nothing
settles the stock purchase. The invention-assignment agreement is signed
separately and is the one that has to be amended, because the schedule of
excluded inventions is not where it binds: it reaches forward to every
invention the faculty member makes at BYU after incorporation, all of it the
university's, and it carries a representation that signing ``does not and will
not breach any agreement or duty which the employee has with anyone else.''
Leaving that schedule blank leaves both of those untouched. So does BYU's
claim, under clause 6.3 of the plan, to a perpetual royalty-free license in
whatever a person not employed by BYU makes in the lab under faculty
supervision.

% ================================================================ failure ==
\parthead{What happens if the plan is not followed}

Two consequences fall on the university rather than on the faculty member.

\begin{itemize}
\item Before any U.S.\ Public Health Service money (the National Institutes of
Health, for example) is spent, BYU must be able to answer public requests about
disclosed financial interests within five business days.
\item Where a financial conflict goes undisclosed, unreviewed, or unmanaged,
or where an approved plan is not followed, the
\href{https://policy.byu.edu/view/financial-conflict-of-interest-in-research-policy}{Financial
Conflict of Interest in Research Policy} calls for a retrospective review
within 120 days of the finding, by a panel of the AAVP, the Research
Administration Office director, the department chair, and the associate dean.
\end{itemize}

% =============================================================== open items
\subhead{Open items}

\begin{itemize}
\item \textbf{For the Technology Transfer Office:} how a forfeited developer
share is handled when a technology has more than one inventor, and whether an
Income Distribution Agreement still circulates once it is forfeited.
\item \textbf{On sequencing, also for that office:} the plan releases openly by
default and sends over only what the published open-source policies do not
cover (clauses 7.1, 7.2). Public release costs foreign patent rights the day it
happens, so the office may want that line drawn elsewhere.
\item \textbf{For the dean and the AAVP:} whether a non-voting board observer
seat is available at all, given that policy addresses only voting seats.
\item \textbf{For Regulatory Accounting:} the published
\href{https://spo.byu.edu/00000190-3687-d777-a591-37af392d0001/rechargecenterpolicyprocedures-pdf}{Recharge
Center Policy} sits outside the policy library entirely (a search of all 195
returns nothing for recharge, service center, or external rate) and still names
a director ``currently (June 2005).'' Confirm the current version before
submitting a rate proposal.
\item \textbf{For Faculty Relations:} the supplemental-compensation cap of one
day a week and the outside-activity cap of four days a month govern different
things but consume the same time. Whether they stack is left open rather than
guessed at.
\item \textbf{For the Research Administration Office:} ``Appendix A, COI
Management Plan Approval and Appeal Process,'' referenced by the Financial
Conflict of Interest in Research Policy but not published.
\item Rules cited from policies behind sign-in at
\href{https://policy.byu.edu}{policy.byu.edu} (campus-office approval, network
commercial use, student employment caps, employee-vendor status) follow an
authenticated review of the full library in July and August 2026.
\end{itemize}

\vfill
{\color{metagray}\footnotesize\noindent
Documents of record: the two \texttt{.tex} sources in
\href{https://github.com/vertical-cloud-lab/byu-vcl}{vertical-cloud-lab/byu-vcl}.
Compiled 27 August 2026.}

\end{document}
